\documentclass[11pt]{article}

\usepackage[utf8]{inputenc}
\usepackage[T1]{fontenc}
\usepackage[french]{babel}
\usepackage{geometry}
\geometry{
    a4paper,
    left=2.5cm,
    right=2.5cm,
    top=2.5cm,
    bottom=2.5cm
}
\usepackage{times}
\renewcommand{\familydefault}{\sfdefault}

\usepackage[svgnames]{xcolor}
\usepackage{graphicx}
\usepackage{listings}

\definecolor{codegreen}{rgb}{0,0.6,0}
\definecolor{codegray}{rgb}{0.5,0.5,0.5}
\definecolor{codepurple}{rgb}{0.58,0,0.82}
\definecolor{backcolour}{rgb}{0.95,0.95,0.92}

\lstdefinestyle{mystyle}{
    backgroundcolor=\color{backcolour},   
    commentstyle=\color{codegreen},
    keywordstyle=\color{magenta},
    numberstyle=\tiny\color{codegray},
    stringstyle=\color{codepurple},
    basicstyle=\ttfamily\footnotesize,
    prebreak=\raisebox{0ex}[0ex][0ex]{\dots},
    breakatwhitespace=true,         
    breaklines=true,                 
    captionpos=b,                    
    keepspaces=true,                 
    numbers=left,                    
    numbersep=5pt,                  
    showspaces=false,                
    showstringspaces=false,
    showtabs=false,                  
    tabsize=2
}

\lstset{style=mystyle}


\newcommand{\ia}[1]{\textcolor{MidnightBlue}{#1}}
\newcommand{\techno}[1]{\texttt{\textcolor{Sienna}{#1}}}
\newcommand{\minio}{\techno{MinIO}}
\newcommand{\docker}{\techno{Docker}}
\newcommand{\openmetadata}{\techno{OpenMetadata}}
\newcommand{\airflow}{\techno{Airflow}}
\newcommand{\boto}{\techno{Boto3}}

\newcommand{\fichier}[1]{\colorbox{LightGrey}{\texttt{\textcolor{FireBrick}{#1}}}}
\newcommand{\bucket}[1]{\colorbox{LightGrey}{\texttt{\textcolor{DodgerBlue}{#1}}}}
\newcommand{\role}[1]{\colorbox{LightGrey}{\texttt{\textcolor{ForestGreen}{#1}}}}


% Penses-bêtes
% % % % %
% FIGURE
% \begin{figure}[p]
%     \centering
%     \includegraphics[width=\textwidth]{../diagrams/DB_schema.drawio.png}
%     \caption{Diagramme du schéma de base de données}
%     \label{fig:db_schema}
% \end{figure}

% CODE
% \begin{lstlisting}[language=bash]
% mc admin policy create local data-engineer-policy policies/DataEngineerReadWrite.json
% mc admin policy attach local data-engineer-policy --user data-engineer
% \end{lstlisting}
% \begin{lstlisting}[language=python]
% mc admin policy create local data-engineer-policy policies/DataEngineerReadWrite.json
% mc admin policy attach local data-engineer-policy --user data-engineer
% \end{lstlisting}

\begin{document}

\title{Modification d'un entrepôt de données}
\author{Gattano César}
\maketitle

\section*{Avant-propos}

Ce rapport a été rédigé en usant de l'intelligence artificielle.
Par soucis de transparence, les parties purement rédigées par l'IA sont
signalées via une couleur de police bleu sombre.
L'auteur certifie que l'ensemble du contenu générée par l'IA a été revu et validé par ses
soins.

Voici un exemple: \ia{Je suis un assistant IA conçu pour vous aider dans la rédaction
    de votre rapport.
    Je génère du texte structuré et cohérent, que vous pouvez ensuite affiner.
    Mon rôle est de vous fournir un cadre de travail, tout en vous laissant le contrôle final.
    Ensemble, nous pouvons créer un document clair et professionnel.
}

En outre, l'auteur indique que l'IA générative a également été utilisée en
auto-complétion durant la rédaction de ce rapport. Cette fonctionnalité ayant
été utilisée tout au long du rapport, aucun indicateur n'a été mis en place
pour souligner cet aspect, si ce n'est cette déclaration.

\tableofcontents

\pagebreak

\section{Introduction}



\end{document}
